\section{Problem Definition}
\label{sec:def}

The input is a preprocessed dataset $X\subseteq[0,1]^d$ of $n$ points. Here, $d$ denotes the number of attributes. We assume that a user is only interested in $d_\text{int}$ among all $d$ attributes, where $d_\text{int}$ is a small number compared to $d$ (e.g., $d_\text{int}=3$ and $d=100$).

\subsection{Terminology}

We refer to a database object as ``\textit{tuple}" and ``\textit{point}" interchangeably throughout the paper. We denote the $i$-th dimensional value of a point $p$ by $p[i]$ where $i \in [1,d]$. For an attribute set $\mathcal{D}$, let $p_{\mathcal{D}}$ denote the projection of $p$ onto $\mathcal{D}$, and let $X_{\mathcal{D}}=\{p_{\mathcal{D}}:p\in X\}$. We remove constant attributes and apply \textit{min-max} normalization to each remaining attribute. Let $a_i$ and $b_i$ denote the minimum and maximum raw values of attribute $i$. A raw value $x[i]$ is mapped to $p[i]=(x[i]-a_i)/(b_i-a_i)$ when larger values are preferred, and to $p[i]=(b_i-x[i])/(b_i-a_i)$ when smaller values are preferred, as with \textit{price}. Thus, every retained attribute ranges from $0$ to $1$, attains both endpoints, and has larger normalized values preferred. Utilities and regret ratios are defined on this normalized representation. As an example, Table~\ref{tab:house} illustrates a database $X=\{p_1, p_2, p_3, p_4, p_5\}$ of $5$ points ($n=5$) where each point is described by $5$ attributes ($d=5$).

\begin{table}[!t]
\centering
\begin{tabular}{|c|c|c|c|c|c|c|}
\hline
\textbf{House} & $D_1$ & $D_2$ & $D_3$ & $D_4$ & $D_5$ & \textbf{Utility} $f(p)$ \\
\hline
$p_1$ & 0.84 & 0.61 & 0.93 & 0.70 & 0.31 & 0.782 \\
\hline
$p_2$ & 0.59 & 0.95 & 0.77 & 0.86 & 0.79 & 0.761 \\
\hline
$p_3$ & 0.69 & 0.84 & 1.00 & 0.99 & 0.55 & 0.820 \\
\hline
$p_4$ & 1.00 & 0.64 & 0.68 & 0.45 & 1.00 & 0.794 \\
\hline
$p_5$ & 0.74 & 1.00 & 0.44 & 1.00 & 0.73 & 0.756 \\
\hline
\end{tabular}
\caption{House dataset with five attributes and utility values for
$u=(0.40,0.35,0.25,0,0)$ ($d_\text{int}=3$).
The attributes $D_1\!\sim\!D_5$ denote \textit{price}, \textit{size}, \textit{commute time}, \textit{age} and \textit{condition score}, respectively.}
\label{tab:house}
\end{table}

Following~\cite{CB00, LC09, NL12, NS10, PW14}, we model user preferences by an unknown nonnegative linear utility function $f(p)=u\cdot p$, where $u\in\mathbb{R}_{\geq0}^{d}$ and $u[i]$ represents the importance of attribute $i$. We use $f$ and its utility vector $u$ interchangeably.

Building on prior studies that individuals consider only a limited number of attributes in high-dimensional decision making tasks~\cite{GG96, GA12, WX22, SC05, TK03, DMW12}, along with our own user survey, we assume that $u$ has exactly $d_\text{int}$ positive entries, where $d_\text{int}\ll d$. From findings from previous studies and a user survey, we adopt an upper bound $d_\text{max}$ of $5$ (i.e., $d_\text{int}\leq d_\text{max}=5$). An attribute (dimension) is said to be \textit{key} if the coefficient of that attribute in the utility vector $u[i]>0$; and \textit{non-key} otherwise.

The candidate attribute set $\mathcal{D}_\text{cand}$ contains all attributes not yet identified as non-key. Here, `candidate' refers to attributes rather than tuples, to avoid confusion.

We define the \textit{regret ratio}\cite{NS10} as follows:

\begin{definition}\cite{NS10}
Given a nonempty set $S \subseteq X$ and a utility vector $u$, the \emph{regret ratio} of $S$ over $X$ w.r.t.\ $u$, denoted by $\mathrm{rr}_X(S, u)$, is defined as
\[
\frac{\max_{p \in X} u \cdot p - \max_{p \in S} u \cdot p}{\max_{p \in X} u \cdot p}
= 1 - \frac{\max_{p \in S} u \cdot p}{\max_{p \in X} u \cdot p}.
\]
\end{definition}

The regret ratio is invariant under positive scaling of $u$. We therefore normalize $u$ so that $\sum_{i=1}^{d}u[i]=1$ ($\lVert u\rVert_1=1$).

For any nonempty subset $S\subseteq X$, the regret ratio lies between $0\%$ and $100\%$. Since each retained attribute attains $1$ and $\lVert u\rVert_1=1$, we have $\max_{p\in X}u\cdot p\geq\max_i u[i]>0$, so the regret ratio is well defined. A lower regret ratio means that the best returned tuple has utility closer to the dataset optimum. A tuple $p^*\in\arg\max_{p\in X}u\cdot p$ is a \textit{favorite tuple}, also called an optimal or maximum-utility tuple.

Consider the example given in Table~\ref{tab:house}. We have the utility vector $u = (0.40, 0.35, 0.25, 0, 0)$. Take the first point $p_1$ for illustration. Its corresponding utility w.r.t.\ $u$ can be calculated as: $f(p_1) = u \cdot p_1 = 0.782$. The utilities of other points are also presented in Table~\ref{tab:house}. The favorite tuple in $X$ is $p_3$, with utility $f(p_3) = 0.820$. If $S = \{p_1, p_2\}$, then the regret ratio of $S$ over $X$ w.r.t.\ $u$ is:
\[
\begin{aligned}
\mathrm{rr}_X(S, u) &= 1 - \frac{\max_{p \in S} u \cdot p}{\max_{p \in X} u \cdot p} = 1-\frac{0.782}{0.820} \approx 4.63\%.
\end{aligned}
\]

Notation is summarized in \iftoggle{TechnicalReport}{Section~\ref{appen:summary} of the \technicalreport}{Appendix~\ref{appen:summary}}.

\subsection{Interaction}

Following interactive regret-minimization frameworks~\cite{NL12, XCL19, WWX21, CW23}, FHDR learns preferences through tuple comparisons. It discovers key attributes and refines their relative utility weights without requiring explicit utility specification.

\paragraph{Format}
Each question presents $s$ tuples from $X$, denoted $\{p_1,p_2,\dots,p_s\}$, described only by attributes in $\mathcal{D}\subseteq\{D_1,D_2,\dots,D_d\}$, where $1\leq\lvert\mathcal{D}\rvert\leq m$. Here, $m$ is the maximum number of attributes displayed per question. The user either selects the most preferred tuple based solely on the displayed attributes or indicates that none of these attributes matters. We call the latter an \textit{irrelevant-attributes response}. This response is distinct from stopping the interaction.

\paragraph{Strong truthfulness.}
Following \cite{XCL19}, an interactive algorithm is \emph{strongly truthful} if it presents only existing objects to the user throughout the interaction. Our algorithm satisfies strong truthfulness because every tuple shown to the user is drawn from the original dataset.

\paragraph{Partial Utility Assumption}\label{par}
We assume a fixed utility vector $u$ and user responses consistent with the partial-utility and attribute-relevance rules below. Users are instructed to treat omitted attributes as identical for each comparison~\cite{chan2024vaccine}, so their common contribution cancels. Following Do et al.~\cite{LA25}, comparisons use the \textit{partial utility} $f_\mathcal{D}(p)=\sum_{i:D_i\in\mathcal{D}}u[i]p[i]$, retaining the original coefficients of the displayed attributes. Users select a tuple maximizing partial utility when at least one displayed attribute is key; otherwise, they give an irrelevant-attributes response.

For the utility vector in Table~\ref{tab:house}, displaying $\mathcal{D}=\{D_1,D_2,D_4\}$ gives $f_\mathcal{D}(p)=0.40p[1]+0.35p[2]$, since $u[4]=0$. The five partial utilities are $0.550$, $0.569$, $0.570$, $0.624$, and $0.646$, respectively, so the user selects $p_5$. If only $D_4$ and $D_5$ are displayed, both weights are zero and the user gives an irrelevant-attributes response.

\paragraph{Information Gained.}
Each response conveys information about the relevance of attributes in $\mathcal{D}$:
\begin{itemize}
    \item A tuple selection indicates that at least one displayed attribute is key.
    \item An irrelevant-attributes response indicates that $u[i]=0$ for every displayed attribute $D_i\in\mathcal{D}$.
\end{itemize}
These responses refine $\mathcal{D}_\text{cand}$. A tuple selection also supplies comparisons on displayed coordinates that can constrain relative utility weights.

\paragraph{Early Termination.}
Users may stop at any time. Before interactive Phase 3, including at zero feedback, \textsc{AttributeSubset} returns a recommendation set using the current candidate attributes. During interactive Phase 3, early stopping returns one recommendation based on the current feasible utility region; exact favorite-tuple retrieval requires sufficient feedback.

\subsection{Objectives}

We consider two objectives based on feedback.

\paragraph{Problem 1: }\label{p1}
With sufficient feedback under the stated response model, our primary objective is to return a favorite tuple $p^*\in\arg\max_{p\in X}u\cdot p$.

\paragraph{Problem 2: }\label{p2}
When interaction ends before interactive Phase 3, we seek a subset $S\subseteq X$ of user-specified size $\lvert S\rvert=K$ with a small regret ratio $\mathrm{rr}_X(S,u)$, without further feedback~\cite{NS10}. This objective includes the zero-feedback case.

